\documentclass[11pt,a4paper]{article}
\usepackage[utf8]{inputenc}
\usepackage{amsmath,amssymb}
\usepackage{geometry}
\geometry{margin=1in}

\title{\textbf{Understanding Transformer Architecture}}
\author{C S Tarun Kumar}
\date{November 20, 2026}

\begin{document}
\maketitle

\section{Overview}
The Transformer architecture (Vaswani et al., 2017) eliminated recurrence in sequence modeling in favor of pure self-attention mechanisms.

\section{Scaled Dot-Product Attention}
The core attention calculation is defined as:
\begin{equation}
\text{Attention}(Q, K, V) = \text{softmax}\left(\frac{QK^T}{\sqrt{d_k}}\right)V
\end{equation}
Where $Q \in \mathbb{R}^{n \times d_k}$ is the Query matrix, $K \in \mathbb{R}^{m \times d_k}$ is the Key matrix, and $V \in \mathbb{R}^{m \times d_v}$ is the Value matrix.

\section{Multi-Head Attention}
Queries, Keys, and Values are linearly projected $h$ times:
\begin{equation}
\text{MultiHead}(Q, K, V) = \text{Concat}(\text{head}_1, \dots, \text{head}_h)W^O
\end{equation}

\section{Positional Encoding}
\begin{align}
PE_{(pos, 2i)} &= \sin\left(\frac{pos}{10000^{2i/d_{\text{model}}}}\right) \\
PE_{(pos, 2i+1)} &= \cos\left(\frac{pos}{10000^{2i/d_{\text{model}}}}\right)
\end{align}

\end{document}
